\documentclass[11pt]{article}
\usepackage[margin=0.86in]{geometry}
\usepackage{amsmath,amssymb,amsthm,mathtools}
\usepackage{lmodern,microtype}
\usepackage[hidelinks]{hyperref}
\setlength{\emergencystretch}{2em}
\title{A six-variable piecewise quadratic function\\
without a polynomial max--min representation}
\author{}
\date{\small Draft for mathematical review}
\begin{document}
\maketitle
\vspace{-1em}

\section{Statement and the whole-space function}

The construction in this note is a continuous piecewise quadratic function on
all of \(\mathbb R^6\), with five polynomial labels including zero. It has no
finite expression in arbitrary real polynomials using pointwise minimum and
maximum. The number, coefficients and degrees of those hypothetical
polynomials are unrestricted. This does not resolve \(\mathrm{PB}(3)\) and
does not assert that six is the smallest counterexample dimension.

Use coordinates \(w=(x,y,z,u,v,r)\) on \(W=\mathbb R^6\). On
\(U=\mathbb R^5\), use the bilinear form

\[
\beta=\operatorname{diag}(1,1,3,-1,-1),
\]

which has signature \((3,2)\). Define the following matrix of three columns:

\[
C(w)=(c_0,c_1,c_2)=
\begin{pmatrix}
x&-v+r/2&-y\\
y&z+u&x\\
z&r/2&0\\
u&-v/2+r&-v/2\\
v&2z+u/2&u/2
\end{pmatrix}.
\tag{1}
\]

Let \(A\) be its first three rows, let \(B\) be its last two rows, and put

\[
Q=C^T\beta C=A^T\operatorname{diag}(1,1,3)A-B^TB,
\qquad D=\det A.
\tag{2}
\]

These definitions use all six independent coordinates. In particular,
\(c_0\) is the projection onto the first five coordinates, and the map
\(w\mapsto C(w)\) is injective.

Direct multiplication gives

\[
Q=\begin{pmatrix}a&b&0\\b&d-a&e\\0&e&d\end{pmatrix},
\quad a=Q_{00},\quad b=Q_{01},\quad d=Q_{22},\quad e=Q_{12}.
\tag{3}
\]

Thus the two identities \(Q_{00}+Q_{11}-Q_{22}=0\) and \(Q_{02}=0\)
hold on the whole space. Define four homogeneous quadratic polynomials by

\[
\begin{aligned}
q_1&=-10b+2d+2e,\\
q_2&=4a+2b-2d-2e,\\
q_3&=-4a+2b+6d-10e,\\
q_4&=-4a+2b+2d+2e.
\end{aligned}
\tag{4}
\]

Write \(h=\min(q_1,q_2,q_3,q_4)\). The fixed function is

\[
f(w)=
\begin{cases}
h(w),&h(w)>0\text{ and }D(w)>0,\\
0,&\text{otherwise}.
\end{cases}
\tag{5}
\]

The positive diagonal coefficient 3 in (2) causes no restriction on the
underlying real quadratic space: multiplying the third row of \(A\) by
\(\sqrt3\) gives the usual form \(A^TA-B^TB\), with the same determinant
sign. All coefficients in (1), (3), and (4) are rational.

\section{Positive Gram margin, continuity, and a finite closed cover}

Here is a matrix identity establishing that the strict inequalities in (4)
imply \(Q\succ0\). Let \(\mathbf1=(1,1,1)^T\), and put

\[
a_i=4e_i-\mathbf1\quad(0\le i<3),\qquad a_3=-\mathbf1,
\]

\[
u_i=(e_i+\mathbf1)/4\quad(0\le i<3),\qquad u_3=0.
\]

For all six pairs \(i<j\), set \(q_{ij}=-a_i^TQa_j\). The identities
\(\sum a_i=0\) and \(\sum u_i a_i^T=I_3\) give

\[
Q=\sum_{i<j}q_{ij}(u_i-u_j)(u_i-u_j)^T.
\tag{6}
\]

For example, this also follows by multiplying out: the matrix
\((a_i^TQa_j)_{i,j=0}^3\) has zero row sums and off-diagonal entries
\(-q_{ij}\), so its quadratic expression is the sum of the corresponding
six squared differences.

The four polynomials in (4) are respectively
\(q_{01},q_{03},q_{12},q_{13}\). On the matrix space (3), the other two satisfy

\[
q_{02}=2q_{01}+6q_{03}+5q_{13},\qquad
q_{23}=q_{01}+2q_{03}+2q_{13}.
\tag{7}
\]

Consequently \(h>0\) makes all six coefficients in (6) strictly positive.
The three vectors \(u_i-u_3\), \(i<3\), span \(\mathbb R^3\), so (6)
implies \(Q\succ0\). If \(A\xi=0\) for a nonzero \(\xi\), (2) instead
gives \(\xi^TQ\xi=-|B\xi|^2\le0\). We have proved

\[
h(w)>0\ \Longrightarrow\ Q(w)\succ0\ \Longrightarrow\ D(w)\ne0.
\tag{8}
\]

The function (5) is therefore continuous on \(\{h>0\}\), since the
determinant sign is locally constant there. It is zero on \(\{h<0\}\), and
continuity at \(h=0\) follows from \(0\le f\le\max(h,0)\).

An explicit closed semialgebraic cover of the whole space consists of

\[
F_0=\{D\le0\}\ \cup\ \bigcup_{i=1}^4\{q_i\le0\},
\]

with polynomial label zero, and

\[
F_i=\{D\ge0,\ q_j\ge0\ (1\le j\le4),\ q_i\le q_j\ (1\le j\le4)\},
\qquad1\le i\le4,
\tag{9}
\]

with label \(q_i\). If \(D=0\) and all four labels are nonnegative, (8)
forces their minimum to be zero. This verifies the labels on every boundary
overlap. The five sets in (9), including \(F_0\), cover all of \(W\).
Furthermore

\[
f(\lambda w)=\lambda^2 f(w)\qquad(\lambda>0).
\tag{10}
\]

For an additional direct check that the Gram-positive set is nonempty, at
\(w=(2,8,0,-1,-7,-9)\) the Gram matrix is

\[
\begin{pmatrix}18&-12&0\\-12&75/2&-3\\0&-3&111/2\end{pmatrix},
\]

whose leading principal minors are \(18,531,58617/2\).

\section{The polynomial Gram-zero curve and all its quadratic equations}

Define the polynomial curve

\[
b(t)=\begin{pmatrix}
t^4-6t^2+1\\
4t^3-4t\\
t^4+2t^2+1\\
2t^4-2\\
4t^3+4t\\
2t^5+4t^3+2t
\end{pmatrix},
\qquad k(t)=(-t,1,1)^T,
\qquad
R_-=(0,1,-1)^T(0,1,-1).
\tag{11}
\]

Write \(d_*(t)=12(1+t^2)^2\). The following are exact polynomial identities:

\[
Q(b(t))=0,\qquad C(b(t))k(t)=0,\qquad
Q(b'(t))=d_*(t)R_-.
\tag{12}
\]

In particular, the column relation is
\(c_1(b(t))=t c_0(b(t))-c_2(b(t))\). Its coefficient of \(c_2\) is
constantly \(-1\); this fact permits the rank-one target \(R_-\) throughout
the curve.

The coefficient matrix of the six entries of \(b(t)\), with rows indexed by
degrees zero through five, is

\[
\begin{pmatrix}
1&0&1&-2&0&0\\
0&-4&0&0&4&2\\
-6&0&2&0&0&0\\
0&4&0&0&4&4\\
1&0&1&2&0&0\\
0&0&0&0&0&2
\end{pmatrix}.
\tag{13}
\]

Its determinant is \(2048=2^{11}\). Thus no nonzero real homogeneous
linear polynomial vanishes identically on the curve. The curve, viewed at a
transcendental ordered parameter, also has no such linear equation.

Let \(I_2\) be the real vector space of homogeneous quadratics \(q\) in the
six original variables satisfying \(q(b(t))\equiv0\). Let

\[
H_2=\{q\in I_2:q(b'(t))\equiv0\}.
\]

The exact spaces needed below are

\[
\dim I_2=10,\qquad
H_2=\operatorname{span}_{\mathbb R}\{Q_{00},Q_{01},Q_{12}+Q_{22}\}.
\tag{14}
\]

These statements concern every real homogeneous quadratic, including all
the additional quadratics outside the Gram span. Here is a finite integer
certificate of (14). Order the 21 monomials \(w_iw_j\) lexicographically by
\(0\le i\le j<6\). Form the \(11\times21\) matrix \(M\) whose column
\((i,j)\) contains the coefficients, in degrees zero through ten, of
\(b_i(t)b_j(t)\). Form the \(9\times21\) matrix \(M'\) from the coefficients,
in degrees zero through eight, of \(b'_i(t)b'_j(t)\), and put
\(N=\binom{M}{M'}\). With all indices beginning at zero:

\begin{itemize}
\item The minor of \(M\) using all rows and columns
  \(0,1,2,3,4,5,6,7,8,9,20\) has determinant \(2^{35}\).
\item The minor of \(N\) using rows \(0,1,\ldots,17\) and columns
  \(0,1,\ldots,16,20\) has determinant \(2^{69}\).
\end{itemize}

The first assertion also follows directly from (13): the coordinate
polynomials span all polynomials of degree at most five, so their products
span all polynomials of degree at most ten. Thus \(\operatorname{rank}M=11\)
and \(\dim I_2=21-11=10\).

By (12), all three displayed Gram polynomials in (14) lie in \(\ker N\).
They are independent: their coefficients on monomials \(x^2,xz,xv\) form

\[
\begin{pmatrix}1&0&0\\0&0&-1\\1&1&0\end{pmatrix},
\]

with determinant one. Hence \(\operatorname{rank}N\le18\), while the second
minor gives the reverse inequality. This proves (14) without assuming that
the four Gram entries generate the full space of quadratic equations.

\section{First and second variations inside the fixed six-dimensional space}

For a homogeneous quadratic matrix map, write \(dQ_b(z)\) for its linear
derivative at \(b\) in direction \(z\). In particular,

\[
Q(b+\varepsilon z)
 =Q(b)+\varepsilon dQ_b(z)+\varepsilon^2Q(z)
\tag{15}
\]

is an exact identity.

Fix arbitrary real \(t\) and \(\rho>0\), and take the base
\(b_*=\rho b(t)\). In this section all derivatives of the curve are
evaluated at that fixed \(t\). Set

\[
n=-\frac{b''}{2\rho d_*},\qquad
j_0=b'''-\frac{6t}{1+t^2}b''.
\tag{16}
\]

Differentiating (12) twice gives
\(dQ_{b(t)}(b'')=-2d_*R_-\). Differentiating a third time gives
\(dQ_{b(t)}(b''')=-3d_*'R_-\). Since
\(3d_*'/(2d_*)=6t/(1+t^2)\), it follows that

\[
dQ_{b_*}(n)=R_-,\qquad dQ_{b_*}(j_0)=0.
\tag{17}
\]

For a symmetric matrix in the four-dimensional space (3), define

\[
\pi(R)=(R_{00},R_{01},R_{22}),\qquad
\ell(R)=k^TRk.
\tag{18}
\]

The maps \(\pi\) and \(\ell\) together identify that matrix space with
\(\mathbb R^4\): explicitly,

\[
\ell(R)=(t^2-1)R_{00}-2tR_{01}+2R_{22}+2R_{12}.
\tag{19}
\]

The equation \(C(b_*)k=0\) implies

\[
\ell(dQ_{b_*}(z))=0\qquad(z\in W).
\tag{20}
\]

Let \(e_x,e_y,e_z\) be the first three standard coordinate vectors of
\(W\), and let \(J_3\) have columns
\(\pi(dQ_{b_*}(e_x)),\pi(dQ_{b_*}(e_y)),\pi(dQ_{b_*}(e_z))\).
An exact derivative calculation gives

\[
\det J_3=12\rho^3(1+t^2)^6>0.
\tag{21}
\]

In particular \(dQ_{b_*}\) has rank exactly three, its image is precisely
\(\ker\ell\), and its three displayed columns give coordinates on that
image. No generic rank hypothesis is being imposed on the chosen base.

The following identities supply the determinant direction and the second
variation missing from (20):

\[
\begin{aligned}
\beta(C(n)k,C(n)k)&=0,\\
\beta(C(n)k,C(j_0)k)&=0,\\
\beta(C(j_0)k,C(j_0)k)&=1728,\\
dD_{b_*}(n)&=0,\\
dD_{b_*}(j_0)&=-48\rho^2(1+t^2)^4.
\end{aligned}
\tag{22}
\]

For the first two identities there is also a direct geometric explanation.
Write \(v_0(t)=c_0(b(t))\). For
\(K=\operatorname{diag}(J,0,J/2)\), with
\(J=\left(\begin{smallmatrix}0&-1\\1&0\end{smallmatrix}\right)\), one has

\[
v_0'=\frac{4t v_0-4Kv_0}{1+t^2},\qquad C(b'')k=2v_0'.
\tag{23}
\]

Thus \(C(n)k\) lies in the isotropic plane spanned by \(v_0,Kv_0\).
These two vectors are independent: the third coordinate of \(v_0\) is
\((1+t^2)^2>0\), the third coordinate of \(Kv_0\) is zero, and
\(x(b(t))^2+y(b(t))^2=(1+t^2)^4>0\) makes \(Kv_0\ne0\).
By (17), \(C(j_0)k\) is perpendicular to that plane. Its positive squared
norm and the two determinant derivatives in (22) are the exact values
given by (1) and (11).

Fix once and for all the positive matrix

\[
R_0=\begin{pmatrix}1&1/4&0\\1/4&1&0\\0&0&2\end{pmatrix}.
\tag{24}
\]

It satisfies (3), and

\[
s=\ell(R_0)=t^2-t/2+3=(t-1/4)^2+47/16>0.
\]

Normalize the determinant direction by

\[
j=-\sqrt{s/1728}\,j_0.
\tag{25}
\]

Equations (17) and (22) now give

\[
\begin{aligned}
dQ_{b_*}(j)&=0,&dD_{b_*}(j)&>0,\\
\ell(Q(n\pm\delta j))&=\delta^2s
&&\text{for every real }\delta.
\end{aligned}
\tag{26}
\]

In particular, the two signs in (26) have exactly the same required second
variation, even when the coefficient \(\delta\) is chosen arbitrarily small.

\section{Corrected paths with exactly the same singularly scaled Gram target}

Fix \(t,\rho\) as above and then fix any \(\delta>0\). Put
\(y_\pm=n\pm\delta j\). The four-dimensional space

\[
E=\operatorname{span}(e_x,e_y,e_z,j)
\]

has the displayed basis: its first three vectors map independently under
\(dQ_{b_*}\), while \(j\) maps to zero and has a nonzero determinant
derivative.

For \(z\in E\), consider the following four real equations:

\[
\Phi_\pm(\varepsilon,z)=
\begin{pmatrix}
\pi\bigl(dQ_{b_*}(z)+\varepsilon[Q(y_\pm+z)-\delta^2R_0]\bigr)\\
\ell\bigl(Q(y_\pm+z)-\delta^2R_0\bigr)
\end{pmatrix}=0.
\tag{27}
\]

At \((\varepsilon,z)=(0,0)\), these equations hold by (26). In the displayed
basis of \(E\), their derivative with respect to \(z\) has the block form

\[
\begin{pmatrix}
J_3&0\\
*&\pm2\delta s
\end{pmatrix}.
\tag{28}
\]

Indeed the derivative of the last component in direction \(j\) is

\[
2\beta(C(n\pm\delta j)k,C(j)k)=\pm2\delta s
\]

by (22) and (25). Thus the determinant of (28) is

\[
\pm24\rho^3(1+t^2)^6\delta s\ne0.
\tag{29}
\]

The ordinary real implicit function theorem supplies smooth solutions
\(z_\pm(\varepsilon)\in E\), with \(z_\pm(0)=0\), on an interval about
zero. These are actual local real paths, not merely formal expansions.
Define points of the fixed whole space by

\[
w_\pm(\varepsilon)=b_*+
\varepsilon\bigl(n\pm\delta j+z_\pm(\varepsilon)\bigr).
\tag{30}
\]

Equations (15), (17), (20), and (27) imply the exact identities

\[
Q(w_\pm(\varepsilon))
 =\varepsilon R_-+\varepsilon^2\delta^2R_0.
\tag{31}
\]

For clarity, the first three components of (27) make the \(\pi\)-values on
the two sides of (31) equal. The fourth component makes their \(\ell\)-values
equal, using (20) and \(\ell(R_-)=0\). Formula (19) then gives equality of
every Gram entry. There is no omitted equation to impose.

Also, (22), (25), and smoothness give

\[
D(w_\pm(\varepsilon))
 =\pm\varepsilon\delta\,dD_{b_*}(j)+O(\varepsilon^2).
\tag{32}
\]

For every sufficiently small positive \(\varepsilon\), their determinant
signs are therefore respectively positive and negative.

The four labels in (4), evaluated on \(R_-\), are \((0,0,16,0)\), and on
\(R_0\) they are \((3/2,1/2,17/2,1/2)\). Hence (31) gives

\[
h(w_+(\varepsilon))=h(w_-(\varepsilon))
  =\tfrac12\varepsilon^2\delta^2>0
\]

for every positive \(\varepsilon\) for which the paths are defined. Thus,
for sufficiently small positive \(\varepsilon\),

\[
f(w_+(\varepsilon))=\tfrac12\varepsilon^2\delta^2>0,
\qquad f(w_-(\varepsilon))=0.
\tag{33}
\]

The normalized target \(Q/\varepsilon\) converges to the rank-one matrix
\(R_-\). The two small eigenvalues of the unnormalized matrix \(Q\) have scale
\(\varepsilon^2\delta^2\). No estimate uniform only near a fixed positive
definite normalized target is used in (27)--(33).

\section{Simultaneous agreement for every finite family of quadratic tests}

Here is the exact finite-family statement:

\begin{quote}
For the fixed whole-space function (5), every finite family of real
polynomials of degree at most two has two points of \(W\) on which all its
members have identical signs, while \(f\) is positive at one point and
zero at the other.
\end{quote}

Fix an arbitrary such family \(\mathcal T\). For a member
\(p=c+\lambda+q\), with homogeneous components of degrees zero, one and two,
consider

\[
p(\rho b(t))=c+\rho\lambda(b(t))+\rho^2q(b(t)).
\tag{34}
\]

If this polynomial in \((\rho,t)\) were identically zero, comparison of
powers of \(\rho\), followed by (13), would give
\(c=0,\lambda=0,q\in I_2\). Thus, for each member of
\(\mathcal T\setminus I_2\), at least one coefficient in \(t\) in (34)
is a nonzero polynomial in \(\rho\). Choose one \(\rho>0\) avoiding the
finitely many roots of these finitely many coefficient polynomials. For
this fixed \(\rho\), every corresponding polynomial (34) is nonzero as a
polynomial in \(t\).

For each \(q\in\mathcal T\cap(I_2\setminus H_2)\), (14) says that
\(q(b'(t))\) is a nonzero polynomial in \(t\). Choose a single real \(t\)
avoiding its roots for all these members, and avoiding the roots of all the
nonzero polynomials (34). Both choices can be made in any prescribed open
intervals of positive \(\rho\) and real \(t\). With \(b_*=\rho b(t)\),
we now have

\[
\begin{aligned}
p(b_*)&\ne0&& (p\in\mathcal T\setminus I_2),\\
q(b')&\ne0&& (q\in\mathcal T\cap(I_2\setminus H_2)).
\end{aligned}
\tag{35}
\]

For a homogeneous quadratic \(q\in I_2\), twice differentiating
\(q(b(t))=0\) gives \(dq_b(b'')=-2q(b')\). The definition of \(n\) in
(16), with the scaling \(b_*=\rho b\), therefore gives

\[
dq_{b_*}(n)=\frac{q(b')}{d_*}.
\tag{36}
\]

For every member in the second line of (35), this is nonzero. Choose one
\(\delta>0\) sufficiently small that

\[
\delta\,|dq_{b_*}(j)|
 <\tfrac12|dq_{b_*}(n)|
\tag{37}
\]

for each of those finitely many members. If the family of constraints is
empty, any positive \(\delta\) may be used. Fix this \(\delta\), and then
construct the two actual paths in Section 5.

For \(q\in\mathcal T\cap(I_2\setminus H_2)\), their expansions are

\[
q(w_\pm(\varepsilon))
 =\varepsilon\bigl(dq_{b_*}(n)\pm\delta dq_{b_*}(j)\bigr)
   +O(\varepsilon^2).
\tag{38}
\]

By (37), the two leading coefficients have the same nonzero sign. Hence the
two values in (38) have the same nonzero sign for sufficiently small
positive \(\varepsilon\). For \(p\in\mathcal T\setminus I_2\), the same
conclusion follows from \(p(b_*)\ne0\) and continuity.

Finally, every \(q\in\mathcal T\cap H_2\) is an actual constant linear
combination of the three Gram entries in (14). Its values on the two paths
are exactly equal by (31), including when that common value is zero. This
handles all extra quadratic equations, not just one selected separator.

Take a single sufficiently small positive \(\varepsilon\) for all these
finitely many conditions, for both implicit solutions, and for (32). All
members of \(\mathcal T\) then have identical signs at the two points,
whereas (33) gives different signs of \(f\). The quantifier order is:

\[
\text{fixed }W,f;\quad
\mathcal T;\quad\rho;\quad t;\quad\delta>0;\quad
\text{the two implicit solutions};\quad\varepsilon>0.
\tag{39}
\]

Neither the function nor its vector space depends on \(\mathcal T\).

\section{Excluding polynomial leaves of arbitrary degrees}

The degree-two tests in Section 6 do not impose a degree bound on the
polynomials in a proposed lattice expression. The following limiting
argument removes all such bounds.

Suppose a positively homogeneous degree-two function \(g\) on a real vector
space has an expression

\[
g=\mathcal L(p_1,\ldots,p_m),
\tag{40}
\]

where \(\mathcal L\) is a finite nonempty expression using minimum and
maximum, and each \(p_i\) is any real polynomial. Write

\[
p_i=c_i+\lambda_i+q_i+\text{terms of degrees at least three},
\]

allowing zero components. At each fixed point \(w\), the limit of
\(\tau^{-2}p_i(\tau w)\), as \(\tau\to0^+\), exists in the extended real
line and is

\[
L_i(w)=
\begin{cases}
\operatorname{sgn}(c_i)\,\infty,&c_i\ne0,\\
\operatorname{sgn}(\lambda_i(w))\,\infty,
     &c_i=0,\ \lambda_i(w)\ne0,\\
q_i(w),&c_i=0,\ \lambda_i(w)=0.
\end{cases}
\tag{41}
\]

Positive scalar multiplication commutes with minimum and maximum. Thus
homogeneity makes (40), after scaling, equal to
\(g(w)=\mathcal L(\tau^{-2}p_1(\tau w),\ldots,\tau^{-2}p_m(\tau w))\)
for every \(\tau>0\). Finite minimum and maximum are continuous on the
extended real line; passing to the limits gives

\[
g(w)=\mathcal L(L_1(w),\ldots,L_m(w)).
\tag{42}
\]

The sign function from the extended real line to the ordered set
\(\{-1,0,1\}\) commutes with minimum and maximum. By (41)--(42), the sign
of \(g(w)\) is therefore determined by the signs of the finite family of
linear and homogeneous quadratic polynomials

\[
\{\lambda_i,q_i:1\le i\le m\},
\tag{43}
\]

together with the fixed constants \(c_i\). This conclusion also covers
vanishing linear or quadratic components and infinite intermediate limits.

Apply the argument to the fixed function \(f\), using (10). Section 6
applied to (43) gives two points on which all those signs agree and the sign
of \(f\) differs. This contradicts (42). Hence (5) has no polynomial lattice
expression with arbitrary finite polynomial leaves.

This is a direct unrestricted-degree nonrepresentation proof. It does not
replace the full separating ideal by a chosen ideal of interface equations,
and it does not claim that a finite family of all possible polynomial tests
can fail to distinguish the two points. Higher-degree test polynomials,
including the determinant itself, may distinguish them; (41)--(43) explain
why that fact cannot produce a lattice representation of this degree-two
homogeneous function.

\section{Exact algebraic checks and scope}

The companion file
\texttt{verify\_certificate6.py} checks (1)--(4), the two
simplex relations, the reconstruction identity, the rational positive Gram
point, all curve and derivative identities, the specified integer minors,
the full quadratic kernel in (14), (21)--(22), and the derivative factors in
(28). Its arithmetic uses symbolic rational numbers and integer polynomial
coefficients; no floating-point rank, sampled sign test, or finite degree
cutoff on polynomial leaves is used.

The continuous whole-space selection, the closed cover, the real implicit
solutions, the finite-family choice of parameters, and the limiting argument
are the proofs in Sections 2, 5, 6, and 7. The result is a six-dimensional
counterexample with homogeneous quadratic labels. It is not a positive or
negative resolution of \(\mathrm{PB}(3)\) or a theorem of global minimality.

\end{document}
